\section{The Fair-Division Legal Argument}
\label{sec:fair-division}

\subsection{The Problem of Algorithmic Justification}

When a redistricting commission submits a plan to a court, it must provide
a justification: why is this plan the best available, and why is the procedure
that produced it neutral?
Algorithmic redistricting provides mathematically precise answers to both
questions, but the precision creates a new challenge: the mathematical
arguments are often inaccessible to non-technical judges and lay jurors.

Consider the \metis{} justification: ``we minimised the total edge-weighted
boundary length of the bisection tree, where edge weights are proportional to
the shared geographic boundary length of adjacent census tracts, using the
multilevel Kernighan-Lin algorithm.''
This is a correct and precise description of what \metis{} does.
But it requires the listener to understand graph partitioning, weighted
adjacency, and multilevel coarsening to evaluate whether it is a good criterion
for compactness.
In \textit{Common Cause v.\ Lewis}, the North Carolina Superior Court
found it necessary to appoint a special master to evaluate the technical
properties of proposed maps --- a symptom of the accessibility gap between
algorithmic precision and legal comprehension.

\subsection{The Moving-Knife Certificate}

The moving-knife protocol provides an alternative justification that is both
mathematically precise and intuitively accessible:

\begin{quote}
\textit{A neutral sweep line was rotated through 180 orientations across the
state.
At each orientation, the line was positioned where it would divide the
state's population equally.
The orientation that produced the most geometrically compact pair of
regions --- as measured by the Reock score, the ratio of each region's
area to the smallest circle that could contain it --- was selected.
Neither party chose the orientation, the line position, or the compactness
criterion: all three were fixed by the algorithm before any district
boundaries were drawn.}
\end{quote}

This description is complete and accurate.
It is also understandable by a non-technical reader: the intuition of a
``neutral sweep'' is familiar from everyday geometric reasoning, and
Reock compactness (``what fraction of the smallest enclosing circle does
the district fill?'') is more concrete than ``minimise total edge-weighted
boundary length.''

\subsection{Rawlsian Procedural Fairness}

The philosophical foundation for the moving-knife certificate is Rawlsian
procedural fairness.
Rawls~\citeyear{rawls1971} argued that a distribution is just if it could
have been chosen by rational agents behind a ``veil of ignorance'' --- that
is, without knowing which position in the distribution they would occupy.
For redistricting, this translates to: a plan is procedurally fair if the
procedure that produced it would have been acceptable to either party before
knowing which districts they would be assigned.

The moving-knife protocol satisfies this criterion.
Before the sweep begins, neither the commission nor any party knows which
half of the state will contain which group of voters --- the partition depends
on the geographic orientation that maximises Reock, which is a function of
tract geometry and population distribution, not of political affiliation.
The two halves of the split are not labelled ``Republican'' or ``Democratic''
during the algorithm's execution; they are labelled only after the plan is
complete.

This stands in contrast to any procedure that selects among candidate plans
based on partisan criteria, which would fail the Rawlsian test because such
a procedure could only be acceptable to the advantaged party.

Puppe and Tasnadi~\citeyear{puppe2026} formalise this argument in the
public-choice setting, proving that a moving-knife redistricting procedure
satisfying anonymity (no party is identified in the algorithm), population
balance (both halves receive equal population), and compactness maximisation
(the most compact orientation is selected) is the unique procedure satisfying
all three axioms simultaneously.
Their axiomatic result provides academic backing for the fair-division
certificate.

\subsection{Comparison to the METIS Argument}

The \metis{} argument is also neutral: minimising edge-weighted boundary
length is a function of geography only, containing no electoral data.
Both arguments are procedurally impartial in the sense that neither
party controls the outcome.
We argue that the \mka{} argument is preferable in three contexts:

\begin{enumerate}
\item \textbf{Jurisdictions where Reock is the operative metric.}
  When a state constitution or court order specifies Reock compactness
  as the required criterion, \mka{} directly satisfies the legal mandate.
  \metis{} does not: it minimises edge cut, not Reock, and cannot certify
  that the submitted plan maximises Reock.

\item \textbf{Courts that prioritise procedural intelligibility.}
  When a court must evaluate the procedure on its own merits --- not just
  the output --- the moving-knife sweep is more accessible than the
  multilevel coarsening of \metis{}.
  The sweep has a direct physical analogy (a ruler rotating across a map)
  that a non-technical judge can visualise; multilevel Kernighan-Lin does not.

\item \textbf{Challenges alleging manipulation.}
  When a plaintiff alleges that the algorithm was configured to favour one
  party, the moving-knife procedure is more resistant to this attack: the
  only configurable parameters are the number of orientations (180 by default)
  and the compactness metric (Reock by default), both of which are publicly
  specified and audit-logged.
  The audit chain records \texttt{optimal\_angle\_deg}, allowing an
  independent auditor to rerun the sweep and confirm that the certified
  angle is correct.
\end{enumerate}

\subsection{Limitations of the Fair-Division Argument}

The moving-knife certificate is not a complete defence in all redistricting
contexts.
Several limitations apply:

\paragraph{VRA majority-minority requirements.}
Section~5 of the Voting Rights Act and state analogues may require that
minority-majority districts be preserved.
\mka{} does not account for racial composition; a plan that maximises Reock
may or may not satisfy VRA requirements.
Practitioners must verify VRA compliance using the platform's VRA analysis
tools (\texttt{bisect analyze --types vra}) separately from the Reock
compactness certification.

\paragraph{Multi-criterion compactness standards.}
Some state constitutions require compactness under multiple metrics.
Florida's Article III, \S 20, has been interpreted to require consideration
of both Reock and Polsby-Popper~\citep{floridaConst}.
A plan certified as Reock-optimal by \mka{} may not be Polsby-Popper-optimal.
In such jurisdictions, the \texttt{--mka-metric polsby} flag or a multi-run
comparison with \cvdgeo{} may be appropriate.

\paragraph{Contiguity and non-convexity.}
For non-convex subgraph geometries (peninsulas, panhandles), a linear sweep
may not produce the globally Reock-optimal split: the best Reock orientation
may depend on the boundary shape of the subgraph in a way that a projection
sweep cannot capture.
\mka{} is a best-available heuristic for non-convex regions, not a global
optimiser.
The audit chain records \texttt{optimal\_angle\_deg} to allow verification
that the sweep was exhaustive at the reported granularity, but does not
certify global optimality.

\paragraph{Bisection tree structure.}
\mka{} applies the moving-knife sweep at each bisection node independently.
The globally Reock-optimal plan may not be achievable by independent local
Reock maximisation at each node, since node-level decisions interact.
This is the same limitation that applies to all recursive bisection algorithms.
